\chapter{Tribes, Regionalism and Ethno-Nationalism}

\section{Regionalism Among Tribals}

\begin{KeyConcept}
\begin{itemize}
\item \textbf{Regionalism} is a form of ethnicity --- an over-attachment to one's region, usually seen as an antithesis of nationalism. Its \textbf{causes (especially among tribals)}:
\item \textbf{Social} --- identity crisis (the feeling that one's culture/language/way of life is under threat; e.g. \textbf{Gorkhaland}, when the West Bengal government imposed Bengali); \textbf{historical} --- the British isolationist policy isolating the North-East (hence the most secessionist movements there); \textbf{political} --- regional identity used for political ends; \textbf{ethnic borders} (e.g. Assam's plain vs hill tribes, with ST-status disputes); \textbf{economic} --- deprivation and 'step-motherly' treatment (Telangana); \textbf{human diversity}; and \textbf{psychological} --- \textbf{primordial loyalties} (Clifford Geertz) to one's ethnic group.
\item \textbf{D. C. Burman} (or A. Burman): regionalism is \textbf{both functional and dysfunctional} --- it maintains diversity but can harm national integrity (or turn secessionist). \textbf{Virginius Xaxa}: for tribals, regional consciousness is a \textbf{sense of identity} --- the demand for autonomy is a way of \textbf{asserting tribal identity}; it is a \textbf{socio-cultural movement against homogenisation}, crystallising 'Adivasi consciousness' at the grassroots due to marginalisation and oppression. Examples: the \textbf{Gonds of Adilabad} (land alienation), the \textbf{Jharkhand Mukti Morcha} (the mouthpiece of tribal rights), and the secessionist \textbf{NSCN}.
\end{itemize}
\end{KeyConcept}

\begin{QuickRevision}
\begin{itemize}
\item Regionalism
\item causes (especially among tribals)
\item Social
\item Gorkhaland
\item historical
\item political
\item ethnic borders
\item economic
\item human diversity
\item psychological
\item primordial loyalties
\item D. C. Burman
\item both functional and dysfunctional
\item Virginius Xaxa
\end{itemize}
\end{QuickRevision}

\section{Ethno-Nationalism}

\begin{KeyConcept}
\begin{itemize}
\item \textbf{Ethno-nationalism} = \textbf{ethnicity + nation}. A \textbf{nation} (Benedict Anderson) is an \textbf{imagined political community that is sovereign and limited} --- an ethnic group that becomes politically aware. \textbf{Ernest Gellner}: a nation is 'building a \textbf{political roof} over one's cultural head'.
\item An \textbf{ethnic group can live without a state} (satisfied with affirmative action), but \textbf{when it claims that its self-realisation depends on mastering the political domain}, it takes steps towards nationhood --- this is \textbf{ethno-nationalism} --- which can express itself by demanding either a \textbf{state within the federal union} or a \textbf{sovereign state}.
\end{itemize}
\end{KeyConcept}

\begin{QuickRevision}
\begin{itemize}
\item Regionalism (tribal): social (identity crisis/Gorkhaland), historical (isolationist), political, ethnic borders (Assam), economic (Telangana), diversity, psychological (primordial loyalties).
\item Burman: functional (diversity) + dysfunctional (secession); Xaxa: identity assertion, anti-homogenisation; Gond, JMM, NSCN.
\item Ethno-nationalism: Anderson (imagined political community, sovereign, limited); Gellner (political roof over cultural head); ethnic group without state vs mastering politics $\rightarrow$ nationhood.
\end{itemize}
\end{QuickRevision}

\section{Religious Change Among Tribes (N. K. Bose's Hindu Mode of Absorption)}

\begin{KeyConcept}
\begin{itemize}
\item \textbf{Religious change among tribes} came from cultural contact with the mainstream: about \textbf{89.4\% of tribals identify as Hindu} and a small minority as Christian. \textbf{Christian missionaries} (funded by the British to pacify tribal revolts) converted the Khasis ($\sim$1813), the Nagas ($\sim$1850) and the Mizos ($\sim$1880) --- offering schools and cooperatives; this was resisted by the \textbf{Birsa Munda/Tana Bhagat movements}.
\item \textbf{N. K. Bose's 'Hindu mode of tribal absorption'}: through contact with plains Hindus, tribes emulate Hindu practices, project a caste identity, and get absorbed. Hindu 'great tradition' absorbed tribal elements --- \textbf{Ganesh} from an elephant totem; \textbf{Kali} as a tribal goddess; myths linking tribes to the epics (Ram meeting the Bhil \textbf{Shabari}, Bhima marrying the tribal woman \textbf{Hidimba}); the \textbf{Bonda} women's 'Sita's curse' legend; and the \textbf{Khasas} claiming descent from the \textbf{Pandavas} to justify polyandry. The \textbf{Hindu Right} later opened gurukuls/Saraswati Shishu Mandirs in tribal areas (Arunachal). \textbf{Christopher von Fürer-Haimendorf} (South Indian Gonds): between 1941 and the 1980s their religion became markedly Hinduised.
\item Later \textbf{Sanskritization} examples: the \textbf{Bauris} of West Bengal (sacred thread, Vaishnavism), the \textbf{Bhumij} of MP (Kshatriya claims --- studied by Surajit Sinha), the \textbf{Pardhan} (Fürer-Haimendorf).
\end{itemize}
\end{KeyConcept}

\begin{QuickRevision}
\begin{itemize}
\item 89.4\% of tribals Hindu; Christian missionaries (Khasis 1813, Nagas 1850, Mizos 1880; Birsa Munda resistance); Hindu Right's gurukuls.
\item N. K. Bose's Hindu mode of tribal absorption; tribal gods $\rightarrow$ great tradition (Ganesh, Kali); Bonda/Sita legend; Khasas (Pandava descent).
\item Sanskritization: Bauris (sacred thread), Bhumij (Kshatriya), Pardhan (Fürer-Haimendorf).
\end{itemize}
\end{QuickRevision}

\section{PYQ Pointers on Tribes}

\begin{KeyConcept}
\begin{itemize}
\item Recurring tribal PYQs and their answers: \textbf{unity and diversity} ($\sim$698 scheduled tribes; linguistic, geographic and economic diversity --- from hunter-gatherers to white-collar); the \textbf{impact of geographical/economic mobility} (the breakdown of the social network $\rightarrow$ community disarticulation; land alienation $\rightarrow$ homelessness/joblessness; from an egalitarian, demonetised, barter society into a hierarchical, monetised, consumerist one --- e.g. the \textbf{Mankidia} tribe displaced by the Similipal tiger reserve, now daily-wage labourers under MGNREGA); \textbf{tribe vs caste} (simple vs complex society, egalitarian vs hierarchical, mechanical vs organic solidarity --- Tönnies's Gemeinschaft/Gesellschaft); \textbf{integration and autonomy} (Elwin's \textbf{isolationist} policy vs Ghurye's \textbf{assimilationist} vs the \textbf{integrationist} policy finally adopted --- the Fifth and Sixth Schedules, tribal sub-plans); \textbf{ethno-nationalism and tribal discontent} (identity crisis, \textbf{jal-jangal-jameen}, economic and political factors); and the \textbf{definitional problem} (evolutionist $\rightarrow$ functionalist $\rightarrow$ the continuum view).
\end{itemize}
\end{KeyConcept}

\begin{QuickRevision}
\begin{itemize}
\item PYQs: unity/diversity (698 tribes); mobility (Mankidia/Similipal); tribe vs caste (Tönnies); integration/autonomy (Elwin vs Ghurye vs integrationist --- Fifth/Sixth Schedule); tribal discontent (identity crisis, jal-jangal-jameen); definitional problem.
\end{itemize}
\end{QuickRevision}
